\documentclass[12pt]{article}

% --- Packages ---
\usepackage[utf8]{inputenc}
\usepackage[T1]{fontenc}
\usepackage[french]{babel}
\usepackage{amsmath, amssymb}
\usepackage{geometry}
\usepackage{tikz}
\usepackage[most]{tcolorbox}
\usepackage{lmodern}
\usepackage{pgfplots}
\pgfplotsset{compat=1.18}
\usepackage{mdframed}
\usepackage{bm}
\usepackage{xcolor}
\usepackage{sectsty}
\usepackage[utf8]{inputenc}
\usepackage[T1]{fontenc}
\usepackage[dvipsnames]{xcolor}
\usepackage{tkz-tab}
\usepackage{multido}
% --- Page ---
\geometry{margin=2cm}

% Footer
\usepackage{fancyhdr}
\pagestyle{fancy}
\cfoot{\thepage}
\rfoot{\textit{G.M.}}
\renewcommand{\headrulewidth}{0pt}
\renewcommand{\footrulewidth}{0pt}

% Pour les normes : 
\newcommand*{\norme}[1]{\left\lVert\overrightarrow{#1}\right\rVert}

\usetikzlibrary{shapes.misc}

\tikzset{cross/.style={cross out, draw=black, minimum size=2*(#1-\pgflinewidth), inner sep=0pt, outer sep=0pt},
%default radius will be 1pt. 
cross/.default={4pt}}

% mettre dans le préambule, après \usepackage{xcolor}
\usepackage{sectsty} % gère facilement la couleur des titres
\setlength{\parindent}{0pt}

\sectionfont{\color{red}\Large\bfseries}     % pour \section*
% \subsectionfont{\color{red}\normalsize\bfseries} % pour \subsection*

\begin{document}

% ============================
%   VECTEURS
% ============================

\begin{center}
    {\LARGE \textbf{Vecteurs}}\\[4mm]
    {\large 2nde générale}
\end{center}

\hrule
\vspace{0.4cm}

% ============================
%   PARTIE 1 : Notions de vecteurs
% ============================

\section*{\textcolor{red}{Partie 1 : Notions de vecteurs}}

\subsection*{\textcolor{red}{1.1 Translation}}

\vspace{5pt}

\begin{mdframed}[
    linewidth=1pt,
    roundcorner=5pt,
    linecolor=ForestGreen
]
{\color{ForestGreen}\textbf{Définition :}}
Une \textbf{translation} fait « glisser » un point ou une figure selon une \underline{direction}, 
un \underline{sens} et une \underline{longueur} donnée, schématisé par \textbf{une flèche}.
\end{mdframed}

\vspace{5pt}

\underline{Exemples :}
% \vspace{5pt}
% \\
\begin{center}
\begin{tikzpicture}[scale=1]

% Axes (facultatif)
% \draw[->] (-0.5,0) -- (4,0) node[right] {$x$};
% \draw[->] (0,-0.5) -- (0,3) node[above] {$y$};

% Points
\coordinate (A) at (1,1);
\coordinate (B) at (3,2);
\coordinate (M) at (-1,2);
\coordinate (M') at (1,3);

% Vecteur AB
\draw[->, ultra thick] (A) -- (B) node[midway, above] {};
\draw[->, ultra thick] (M) -- (M') node[midway, above] {};

% Points visibles
\fill (B) circle (1pt) node[right] {$B$};
\fill (A) circle (1pt) node[left] {$A$};
\fill (M) circle (1pt) node[left] {$M$};
\fill (M') circle (1pt) node[right] {$M'$};

\end{tikzpicture}
\end{center}

\underline{Vocabulaire :} On dit alors que $M'$ est l'image de $M$ par la translation de vecteur
$\overrightarrow{AB}$.

\vspace{10pt}

% ------------------------
% Définition et propriétés
% ------------------------

\subsection*{\textcolor{red}{1.2 Définition et propriétés des vecteurs}}

\vspace{5pt}

\begin{mdframed}[
    linewidth=1pt,
    roundcorner=5pt,
    linecolor=ForestGreen
]
{\color{ForestGreen}\textbf{Définition :}}
Soient $A$ et $B$ deux points distincts. Le vecteur $\overrightarrow{AB}$ se définit par :
\begin{itemize}
    \item une direction : la droite $(AB)$,
    \item un sens : du point $A$ vers le point $B$,
    \item une longueur : la \textbf{norme} du vecteur $\overrightarrow{AB}$ notée $\norme{AB}$
\end{itemize}
\end{mdframed}

\vspace{5pt}

{\color{red}\underline{Attention :} } Ne pas confondre \textbf{direction} et \textbf{sens} !
\vspace{10pt}
\\La droite $(AB)$ définit \_\_\_\_\_\_\_\_\_\_\_\_ , tandis que la flèche définit 
\_\_\_\_\_\_\_\_\_\_\_\_ : \textbf{de $\bm{A}$ vers $\bm{B}$} ($\overrightarrow{AB}$)

\vspace{10pt}

\underline{Remarques :}
\vspace{5pt}
\begin{itemize}
    \item Si les points $A$ et $B$ sont confondus, alors le vecteur $\overrightarrow{AB}$ devient le
    vecteur $\overrightarrow{AA}$ qui est le 
    \vspace{5pt}
    \\\textbf{vecteur nul}. $\overrightarrow{AA} = $ 
    \_\_\_\_.
\end{itemize}

% \vspace{10pt}
\newpage

\begin{mdframed}[
    linewidth=1pt,
    roundcorner=5pt,
    linecolor=RoyalBlue,
    innertopmargin=10pt,
    innerbottommargin=10pt
]
{\color{RoyalBlue}\textbf{\underline{Exercice d'application n°1 :} Construire l'image d'une figure par une 
translation.}}
\vspace{10pt}

\begin{minipage}{.55\textwidth}%
    \begin{tikzpicture}[scale=0.5]

    % Grille seule
    \draw[step=1cm, black, black, thin] (0,0) grid (17,10);

    % Points
    \coordinate (A) at (1,7);
    \coordinate (B) at (5,8);
    \coordinate (C) at (3,6);
    \coordinate (D) at (6,2);
    \coordinate (F) at (9,4);
    \coordinate (E) at (9,1);

    % Vecteur
    \draw[->, ultra thick] (A) -- (B);
    \draw[-, thick] (C) -- (D);
    \draw[-, thick] (D) -- (E);
    \draw[-, thick] (E) -- (F);
    \draw[-, thick] (C) -- (F);

    \fill (A) circle (1pt) node[left] {};
    \fill (B) circle (1pt) node[left] {};
    \fill (C) circle (1pt) node[left] {};
    \fill (D) circle (1pt) node[left] {};
    \fill (E) circle (1pt) node[left] {};
    \fill (F) circle (1pt) node[left] {};

    % Noms des points
    \node[below left] at (A) {$A$};
    \node[above right] at (B) {$B$};
    \node[below left] at (C) {$C$};
    \node[below left] at (D) {$D$};
    \node[below left] at (E) {$E$};
    \node[below left] at (F) {$F$};

    \end{tikzpicture}
\end{minipage}%
\hfill
\begin{minipage}{.45\textwidth}%
    Soit la translation définie par le vecteur $\overrightarrow{AB}$. Construire l'image 
    $C'$, $D'$, $E'$ et $F'$ du quadrilatère $CDEF$ par cette translation.
\end{minipage}%

\vspace{20pt}

{\color{RoyalBlue}\textbf{\underline{Exercice d'application n°2 :} Construire l'image d'un point par une 
translation.}}

\vspace{10pt}

\begin{minipage}{.35\textwidth}%
    \begin{center}
        \begin{tikzpicture}[scale=1]

        % Grille seule
        \draw[step=1cm, black, black, thin] (0,0) grid (3,2);

        % Points
        \coordinate (A) at (0,0);
        \coordinate (B) at (1,0);
        \coordinate (C) at (2,0);
        \coordinate (D) at (3,0);
        \coordinate (E) at (3,1);
        \coordinate (F) at (2,1);
        \coordinate (G) at (1,1);
        \coordinate (H) at (0,1);
        \coordinate (I) at (0,2);
        \coordinate (J) at (1,2);
        \coordinate (K) at (2,2);
        \coordinate (L) at (3,2);

        \fill (A) circle (1pt) node[left] {};
        \fill (B) circle (1pt) node[left] {};
        \fill (C) circle (1pt) node[left] {};
        \fill (D) circle (1pt) node[left] {};
        \fill (E) circle (1pt) node[left] {};
        \fill (F) circle (1pt) node[left] {};
        \fill (G) circle (1pt) node[left] {};
        \fill (H) circle (1pt) node[left] {};
        \fill (I) circle (1pt) node[left] {};
        \fill (J) circle (1pt) node[left] {};
        \fill (K) circle (1pt) node[left] {};
        \fill (L) circle (1pt) node[left] {};

        % Noms des points
        \node[below left] at (A) {$A$};
        \node[below left] at (B) {$B$};
        \node[below left] at (C) {$C$};
        \node[below left] at (D) {$D$};
        \node[below left] at (E) {$E$};
        \node[below left] at (F) {$F$};
        \node[below left] at (G) {$G$};
        \node[below left] at (H) {$H$};
        \node[below left] at (I) {$I$};
        \node[below left] at (J) {$J$};
        \node[below left] at (K) {$K$};
        \node[below left] at (L) {$L$};

        \end{tikzpicture}
    \end{center}
\end{minipage}%
\hfill
\begin{minipage}{.65\textwidth}%
    6 carrés sont juxtaposés.
    \\Donner l'image : 
    \begin{enumerate}
        \item de $B$ par la translation de vecteur $\overrightarrow{GF}$
        \item de $D$ par la translation de vecteur $\overrightarrow{KI}$
        \item de $G$ par la translation de vecteur $\overrightarrow{CE}$
        \item de $E$ par la translation de vecteur $\overrightarrow{KH}$
    \end{enumerate}
\end{minipage}%

\vspace{20pt}

{\color{RoyalBlue}\textbf{\underline{Exercice d'application n°3 :} Placer un point par une translation.}}

\vspace{10pt}

\begin{minipage}{.30\textwidth}%
    \begin{center}
        \begin{tikzpicture}[scale=0.7]

        % Grille seule
        \draw[step=1cm, black, black, thin] (0,0) grid (6,4);

        % Points
        \coordinate (F) at (1,2);
        \coordinate (G) at (4,3);
        \coordinate (E) at (2,1);

        \fill (F) circle (2pt) node[left] {};
        \fill (G) circle (2pt) node[left] {};
        \fill (E) circle (2pt) node[left] {};

        % Noms des points
        \node[below left] at (F) {$F$};
        \node[below left] at (G) {$G$};
        \node[below left] at (E) {$E$};

        \end{tikzpicture}
        \\$T$ est l'image de $E$ par la translation qui à $F$ associe $G$.
    \end{center}
\end{minipage}%
\hfill
\begin{minipage}{.30\textwidth}%
    \begin{center}
        \begin{tikzpicture}[scale=0.7]

        % Grille seule
        \draw[step=1cm, black, black, thin] (0,0) grid (6,4);

        % Points
        \coordinate (D) at (3,3);
        \coordinate (B) at (6,2);
        \coordinate (C) at (2,1);

        \fill (D) circle (2pt) node[left] {};
        \fill (B) circle (2pt) node[left] {};
        \fill (C) circle (2pt) node[left] {};

        % Noms des points
        \node[below left] at (D) {$D$};
        \node[below left] at (B) {$B$};
        \node[below left] at (C) {$C$};

        \end{tikzpicture}
        \\$K$ est l'image de $C$ par la translation de vecteur $\overrightarrow{DB}$.
    \end{center}
\end{minipage}%
\hfill
\begin{minipage}{.30\textwidth}%
    \begin{center}
        \begin{tikzpicture}[scale=0.7]

        % Grille seule
        \draw[step=1cm, black, black, thin] (0,0) grid (6,4);

        % Points
        \coordinate (A) at (3,3);
        \coordinate (B) at (1,2);
        \coordinate (C) at (4,1);

        \fill (A) circle (2pt) node[left] {};
        \fill (B) circle (2pt) node[left] {};
        \fill (C) circle (2pt) node[left] {};

        % Noms des points
        \node[below left] at (A) {$A$};
        \node[below left] at (B) {$B$};
        \node[below left] at (C) {$C$};

        \end{tikzpicture}
        \\$P$ est l'image de $A$ par la translation de vecteur $\overrightarrow{BC}$.
    \end{center}
\end{minipage}%

\end{mdframed}

\vspace{10pt}

% ------------------------
% Vecteurs égaux
% ------------------------

\subsection*{\textcolor{red}{1.3 Vecteurs égaux}}

\vspace{5pt}

\begin{mdframed}[
    linewidth=1pt,
    roundcorner=5pt,
    linecolor=ForestGreen
]
{\color{ForestGreen}\textbf{Définition :}}
Des vecteurs sont \textbf{égaux} lorsqu'ils ont même direction, même sens et même longueur. Ils définissent 
alors la même translation.
\end{mdframed}

% \vspace{5pt}
\newpage

\underline{Exemple :}
\vspace{5pt}
\\
\begin{minipage}{.45\textwidth}%
Les vecteurs $\overrightarrow{AB}$ et $\overrightarrow{CD}$ sont égaux.
\\On note : $\overrightarrow{AB} = \overrightarrow{CD}$.
\end{minipage}%
\hfill
\begin{minipage}{.55\textwidth}%
\begin{center}
\begin{tikzpicture}[scale=1]

% Axes (facultatif)
% \draw[->] (-0.5,0) -- (4,0) node[right] {$x$};
% \draw[->] (0,-0.5) -- (0,3) node[above] {$y$};

\draw[step=1cm, black, thin] (-2.5,-0.5) grid (4.5,4.5);

% Points
\coordinate (A) at (1,1);
\coordinate (B) at (3,2);
\coordinate (C) at (-1,2);
\coordinate (D) at (1,3);

% Vecteur AB
\draw[->, ultra thick] (A) -- (B) node[midway, above, xshift=-3pt, yshift=1pt] {$\overrightarrow{AB}$};
\draw[->, ultra thick] (C) -- (D) node[midway, above, xshift=-3pt, yshift=1pt] {$\overrightarrow{CD}$};

% Points visibles
\fill (B) circle (1pt) node[right] {$B$};
\fill (A) circle (1pt) node[left] {$A$};
\fill (C) circle (1pt) node[left] {$C$};
\fill (D) circle (1pt) node[right] {$D$};

\end{tikzpicture}
\end{center}
\end{minipage}%

\vspace{20pt}

\underline{Vocabulaire :}
\vspace{5pt}
\\Comme tous les vecteurs qui définissent une même translation sont égaux, le vecteur choisi pour définir la
translation est un \textbf{représentant} de tous les vecteurs égaux.
\vspace{5pt}
\\Il est souvent noté \_\_\_\_\_\_\_.


\vspace{20pt}

\underline{Exemples :}
\vspace{5pt}
\\
\begin{minipage}[t]{.45\textwidth}%
    \vspace{0pt}
    \begin{tikzpicture}[scale=1]

    % Grille seule
    \draw[step=1cm, black, black, thin] (0.5,0.5) grid (8.5,7.5);

    % Points
    \coordinate (A) at (2,4);
    \coordinate (B) at (3,2);
    \coordinate (H) at (1,3);
    \coordinate (F) at (5,2);
    \coordinate (G) at (6,4);
    \coordinate (E) at (3,6);
    \coordinate (I) at (6,6);
    \coordinate (J) at (7,4);
    % \coordinate (G) at (6,4);

    % Vecteur AB
    \draw[->, ultra thick] (A) -- (B) node[midway, above, xshift=-3pt, yshift=1pt] {};
    \draw[->, ultra thick] (I) -- (J) node[midway, above, xshift=5pt, yshift=0pt] {$\overrightarrow{u}$};

    \fill (A) circle (1pt) node[left] {};
    \fill (B) circle (1pt) node[left] {};
    \fill (H) circle (2pt) node[left] {};
    \fill (F) circle (2pt) node[left] {};
    \fill (G) circle (2pt) node[left] {};
    \fill (E) circle (2pt) node[left] {};

    % Noms des points
    \node[below left] at (A) {$A$};
    \node[below left] at (B) {$B$};
    \node[below left] at (H) {$H$};
    \node[below left] at (G) {$G$};
    \node[below left] at (F) {$F$};
    \node[below left] at (E) {$E$};

    \end{tikzpicture}
\end{minipage}%
\hfill
\begin{minipage}[t]{.50\textwidth}%
    \vspace{10pt}
    \underline{Sur la figure ci-contre :}
    \vspace{5pt}
    \\Le vecteur $\overrightarrow{AB}$ est le représentant d'origine A du vecteur $\overrightarrow{u}$.
    \vspace{5pt}
    \begin{enumerate}
        \item Construire le vecteur représentant de $\overrightarrow{u}$ d'extrémité $G$.
        \item Construire le point $E'$ tel que $\overrightarrow{EE'} = \overrightarrow{u}$.
        \item Construire le point $H'$ tel que $\overrightarrow{H'H} = \overrightarrow{u}$.
        \item Construire le point $F'$ tel que $\overrightarrow{F'F}$ soit égal à l'opposé de 
        $\overrightarrow{u}$.
    \end{enumerate}
\end{minipage}%

\vspace{20pt}


% --------------------------
% ------ GRAPHIQUE 1 -------
% --------------------------
\begin{center}
    \begin{minipage}[t]{.5\textwidth}%
        \centering
        \vspace{0pt}
        \begin{tikzpicture}[scale=1]

        % Grille seule
        \draw[step=1cm, black, black, thin] (0.5,0.5) grid (4.5,3.5);

        % Points
        \coordinate (A) at (1,1);
        \coordinate (B) at (2,3);
        \coordinate (H) at (2,1);
        \coordinate (F) at (4,1);
        % \coordinate (G) at (6,4);

        % Vecteur AB
        \draw[->, ultra thick] (A) -- (B) node[midway, above, xshift=-3pt, yshift=1pt] {$\overrightarrow{u}$};
        \draw[->, ultra thick] (H) -- (F) node[midway, above, xshift=5pt, yshift=0pt] {$\overrightarrow{v}$};

        % \fill (A) circle (1pt) node[left] {};
        % \fill (B) circle (1pt) node[left] {};
        % \fill (H) circle (2pt) node[left] {};
        % \fill (F) circle (2pt) node[left] {};

        % Noms des points
        % \node[below left] at (A) {$A$};
        % \node[below left] at (B) {$B$};
        % \node[below left] at (H) {$H$};
        % \node[below left] at (F) {$F$};

        \end{tikzpicture}
    \end{minipage}%
    \hfill
    \begin{minipage}[t]{.5\textwidth}%
        \vspace{5pt}
        \begin{mdframed}
            Les vecteurs $\overrightarrow{u}$ et $\overrightarrow{v}$ n'ont :
            \begin{itemize}
                \item pas le même sens
                \item pas la même direction
                \item pas la même norme
            \end{itemize}
        \end{mdframed}
    \end{minipage}%
\end{center}

\vspace{20pt}

% --------------------------
% ------ GRAPHIQUE 2 -------
% --------------------------
\begin{minipage}[t]{.35\textwidth}%
    \centering
    \vspace{10pt}
    \begin{tikzpicture}[scale=0.8]

    % Grille seule
    \draw[step=1cm, black, black, thin] (0.5,-0.5) grid (4.5,3.5);

    % Points
    \coordinate (A) at (1,1);
    \coordinate (B) at (2,3);
    \coordinate (H) at (2,1);
    \coordinate (F) at (4,0);
    % \coordinate (G) at (6,4);

    % Vecteur AB
    \draw[->, ultra thick] (A) -- (B) node[midway, above, xshift=-3pt, yshift=1pt] {$\overrightarrow{u}$};
    \draw[->, ultra thick] (H) -- (F) node[midway, above, xshift=5pt, yshift=0pt] {$\overrightarrow{v}$};

    \end{tikzpicture}
\end{minipage}%
\hfill
\begin{minipage}[t]{.6\textwidth}%
    \vspace{5pt}
    \begin{mdframed}
        Les vecteurs $\overrightarrow{u}$ et $\overrightarrow{v}$ n'ont :
        \begin{itemize}
            \vspace{7pt}
            \item pas \_\_\_\_\_\_\_\_\_\_\_\_\_\_\_\_\_\_\_\_\_\_\_\_\_
            \vspace{7pt}
            \item pas \_\_\_\_\_\_\_\_\_\_\_\_\_\_\_\_\_\_\_\_\_\_\_\_\_
        \end{itemize}
        \vspace{7pt}
        En revanche, les vecteurs $\overrightarrow{u}$ et $\overrightarrow{v}$ ont 
        \vspace{7pt}
        \\
        \_\_\_\_\_\_\_\_\_\_\_\_\_\_\_\_\_\_\_\_\_\_\_\_\_.
    \end{mdframed}
\end{minipage}%

\vspace{30pt}

% --------------------------
% ------ GRAPHIQUE 3 -------
% --------------------------
\begin{minipage}[t]{.35\textwidth}%
    \centering
    \vspace{0pt}
    \begin{tikzpicture}[scale=0.8]

    % Grille seule
    \draw[step=1cm, black, black, thin] (0.5,-1.5) grid (4.5,3.5);

    % Points
    \coordinate (A) at (1,1);
    \coordinate (B) at (2,3);
    \coordinate (H) at (4,3);
    \coordinate (F) at (2,-1);
    % \coordinate (G) at (6,4);

    % Vecteur AB
    \draw[->, ultra thick] (A) -- (B) node[midway, above, xshift=-3pt, yshift=1pt] {$\overrightarrow{u}$};
    \draw[->, ultra thick] (H) -- (F) node[midway, above, xshift=-5pt, yshift=0pt] {$\overrightarrow{v}$};

    \end{tikzpicture}
\end{minipage}%
\hfill
\begin{minipage}[t]{.6\textwidth}%
    \vspace{0pt}
    \begin{mdframed}
        Les vecteurs $\overrightarrow{u}$ et $\overrightarrow{v}$ n'ont pas 
        \vspace{7pt}
        \\
        \_\_\_\_\_\_\_\_\_\_\_\_\_\_\_\_\_\_\_\_\_\_\_\_\_.
        \vspace{7pt}
        \\
        Les vecteurs $\overrightarrow{u}$ et $\overrightarrow{v}$ ont 
        \vspace{7pt}
        \\
        \_\_\_\_\_\_\_\_\_\_\_\_\_\_\_\_\_\_\_\_\_\_\_\_\_
        \vspace{10pt}
        \\
        \_\_\_\_\_\_\_\_\_\_\_\_\_\_\_\_\_\_\_\_\_\_\_\_\_.
    \end{mdframed}
\end{minipage}%

\vspace{30pt}

% --------------------------
% ------ GRAPHIQUE 4 -------
% --------------------------
\begin{minipage}[t]{.35\textwidth}%
    \centering
    \vspace{0pt}
    \begin{tikzpicture}[scale=0.8]

    % Grille seule
    \draw[step=1cm, black, black, thin] (0.5,-1.5) grid (4.5,3.5);

    % Points
    \coordinate (A) at (1,0);
    \coordinate (B) at (2,2);
    \coordinate (H) at (4,3);
    \coordinate (F) at (2,-1);
    % \coordinate (G) at (6,4);

    % Vecteur AB
    \draw[->, ultra thick] (A) -- (B) node[midway, above, xshift=-3pt, yshift=1pt] {$\overrightarrow{u}$};
    \draw[->, ultra thick] (F) -- (H) node[midway, above, xshift=-5pt, yshift=0pt] {$\overrightarrow{v}$};

    \end{tikzpicture}
\end{minipage}%
\hfill
\begin{minipage}[t]{.6\textwidth}%
    \vspace{0pt}
    \begin{mdframed}
        Les vecteurs $\overrightarrow{u}$ et $\overrightarrow{v}$ n'ont pas 
        \vspace{7pt}
        \\
        \_\_\_\_\_\_\_\_\_\_\_\_\_\_\_\_\_\_\_\_\_\_\_\_\_.
        \vspace{7pt}
        \\
        Les vecteurs $\overrightarrow{u}$ et $\overrightarrow{v}$ ont 
        \vspace{7pt}
        \\
        \_\_\_\_\_\_\_\_\_\_\_\_\_\_\_\_\_\_\_\_\_\_\_\_\_
        \vspace{10pt}
        \\
        \_\_\_\_\_\_\_\_\_\_\_\_\_\_\_\_\_\_\_\_\_\_\_\_\_.
    \end{mdframed}
\end{minipage}%

\vspace{30pt}

% --------------------------
% ------ GRAPHIQUE 5 -------
% --------------------------
\begin{minipage}[t]{.35\textwidth}%
    \centering
    \vspace{0pt}
    \begin{tikzpicture}[scale=0.8]

    % Grille seule
    \draw[step=1cm, black, black, thin] (0.5,-1.5) grid (4.5,3.5);

    % Points
    \coordinate (A) at (1,0);
    \coordinate (B) at (2,2);
    \coordinate (H) at (3,1);
    \coordinate (F) at (2,-1);
    % \coordinate (G) at (6,4);

    % Vecteur AB
    \draw[->, ultra thick] (A) -- (B) node[midway, above, xshift=-3pt, yshift=1pt] {$\overrightarrow{u}$};
    \draw[->, ultra thick] (F) -- (H) node[midway, above, xshift=-5pt, yshift=0pt] {$\overrightarrow{v}$};

    \end{tikzpicture}
\end{minipage}%
\hfill
\begin{minipage}[t]{.6\textwidth}%
    \vspace{0pt}
    \begin{mdframed}
        \vspace{20pt}
        \_\_\_\_\_\_\_\_\_\_\_\_\_\_\_\_\_\_\_\_\_\_\_\_\_\_\_\_\_\_\_
        \vspace{10pt}
        \\\_\_\_\_\_\_\_\_\_\_\_\_\_\_\_\_\_\_\_\_\_\_\_\_\_\_\_\_\_\_\_
        \vspace{10pt}
        \\\_\_\_\_\_\_\_\_\_\_\_\_\_\_\_\_\_\_\_\_\_\_\_\_\_\_\_\_\_\_\_
        \vspace{10pt}
        \\\_\_\_\_\_\_\_\_\_\_\_\_\_\_\_\_\_\_\_\_\_\_\_\_\_\_\_\_\_\_\_
        \vspace{10pt}
        \\\_\_\_\_\_\_\_\_\_\_\_\_\_\_\_\_\_\_\_\_\_\_\_\_\_\_\_\_\_\_\_
    \end{mdframed}
\end{minipage}%

\newpage

\begin{mdframed}[
    linewidth=1pt,
    roundcorner=5pt,
    linecolor=Red,
    innertopmargin=10pt,
    innerbottommargin=10pt
]
{\color{red}
\underline{Propriété du parallélogramme:}
\vspace{5pt}
\\Dire que $\overrightarrow{AB}$ et $\overrightarrow{CD}$ sont égaux revient à dire que le quadrilatère
$AB\bm{DC}$ est un parallélogramme, éventuellement aplati.
}
\\
\begin{center}
    \begin{tikzpicture}[scale=1]

    % Grille (optionnelle, pédagogique)
    \draw[step=1cm, black, thin] (-0.5,-0.5) grid (6.5,4.5);

    % Points
    \coordinate (A) at (1,1);
    \coordinate (B) at (4,2);
    \coordinate (C) at (2,3);
    \coordinate (D) at (5,4);

    % Parallélogramme (côtés)
    % \draw[thick] (A) -- (B) -- (D) -- (C) -- cycle;
    \draw[dashed, thick] (A) -- (C);
    \draw[dashed, thick] (D) -- (B);

    % Vecteurs égaux
    \draw[->, ultra thick, blue] (A) -- (B)
        node[midway, below] {$\overrightarrow{AB}$};
    \draw[->, ultra thick, blue] (C) -- (D)
        node[midway, above] {$\overrightarrow{CD}$};

    % Diagonales en pointillés
    % \draw[dashed, thick] (A) -- (D);
    % \draw[dashed, thick] (B) -- (C);

    % Points visibles
    \fill (A) circle (1pt) node[below left] {$A$};
    \fill (B) circle (1pt) node[below right] {$B$};
    \fill (C) circle (1pt) node[above left] {$C$};
    \fill (D) circle (1pt) node[above right] {$D$};

    \end{tikzpicture}
\end{center}
\end{mdframed}

\vspace{10pt}

\begin{mdframed}[
    linewidth=1pt,
    roundcorner=5pt,
    linecolor=Red,
    innertopmargin=10pt,
    innerbottommargin=10pt
]
{\color{red}
\underline{Propriété du milieu :}
\vspace{5pt}
\\Dire que $B$ est le milieu du segment $\left[AC\right]$ revient à dire que $\overrightarrow{AB}$ et 
$\overrightarrow{BC}$ sont égaux.
}
\\
\begin{center}
    \begin{tikzpicture}[scale=1]

    % Grille (optionnelle, pédagogique)
    % \draw[step=1cm, black, thin] (-0.5,0.5) grid (5.5,2.5);

    % Points
    \coordinate (A) at (1,1);
    \coordinate (B) at (2.5,1);
    \coordinate (C) at (4,1);

    % Vecteurs égaux
    \draw[->, ultra thick, blue] (A) -- (B)
        node[midway, below] {$\overrightarrow{AB}$};
    \draw[->, ultra thick, blue] (B) -- (C)
        node[midway, below] {$\overrightarrow{BC}$};

    % Points visibles
    \fill (A) circle (1pt) node[above] {$A$};
    \fill (B) circle (1pt) node[above] {$B$};
    \fill (C) circle (1pt) node[above] {$C$};

    \end{tikzpicture}
\end{center}
\end{mdframed}

\vspace{10pt}

% EXERCICE D'APPLICATION

\begin{mdframed}[
    linewidth=1pt,
    roundcorner=5pt,
    linecolor=RoyalBlue,
    innertopmargin=10pt,
    innerbottommargin=10pt
]
{\color{RoyalBlue}\textbf{\underline{Exercice d'application n°4 :}}
\vspace{5pt}
$ABCD$ et $ABDE$ sont deux parallélogramme.
\begin{enumerate}
    \item Réaliser une figure.
    \item Démontrer que $D$ est le milieu de $\left[EC\right]$.
    \item $F$ est le symétrique de $B$ par rapport à $D$.
    \\Démontrer que $DFEA$ est un parallélogramme.
\end{enumerate}
}

\newpage

% EXERCICE N°5
{\color{RoyalBlue}\textbf{\underline{Exercice d'application n°5 :}}
\vspace{5pt}
Dire si les égalités suivantes sont vraies ou fausses.
}

\begin{center}
\begin{minipage}{.24\textwidth}%
\centering
\textbf{1}

\vspace{10pt}

\begin{tikzpicture}[scale=0.8]
\draw[step=1cm, black] (0,0) grid (4,4);

\coordinate (A) at (1,0);
\coordinate (B) at (3,1);
\coordinate (M) at (1,3);
\coordinate (N) at (3,2);

\draw[->, very thick] (A)--(B) node[midway, below right] {};
\draw[->, very thick] (M)--(N) node[midway, below left] {};

\fill (A) circle (1pt) node[below left] {$A$};
\fill (B) circle (1pt) node[above right] {$B$};
\fill (M) circle (1pt) node[below left] {$M$};
\fill (N) circle (1pt) node[below left] {$N$};
\end{tikzpicture}

\vspace{10pt}

$\overrightarrow{AB} = \overrightarrow{MN}$
\end{minipage}%
\hfill
\begin{minipage}{.24\textwidth}%
\centering
\textbf{2}

\vspace{10pt}

\begin{tikzpicture}[scale=0.8]
\draw[step=1cm, black] (0,0) grid (4,4);

\coordinate (A) at (1,2);
\coordinate (B) at (3,2);
\coordinate (M) at (3,3);
\coordinate (N) at (1,3);

\draw[->, very thick] (A)--(B) node[midway, below] {};
\draw[->, very thick] (M)--(N) node[midway, above] {};

\fill (A) circle (1pt) node[below left] {$A$};
\fill (B) circle (1pt) node[above right] {$B$};
\fill (M) circle (1pt) node[above right] {$M$};
\fill (N) circle (1pt) node[below left] {$N$};
\end{tikzpicture}

\vspace{10pt}

$\overrightarrow{AB} = \overrightarrow{MN}$
\end{minipage}%
\hfill
\begin{minipage}{.24\textwidth}%
\centering
\textbf{3}

\vspace{10pt}

\begin{tikzpicture}[scale=0.8]
\draw[step=1cm, black] (0,0) grid (4,4);

\coordinate (A) at (1,1);
\coordinate (B) at (3,2);
\coordinate (M) at (1,3);
\coordinate (N) at (3,4);

\draw[->, very thick] (A)--(B) node[midway, below] {};
\draw[->, very thick] (M)--(N) node[midway, above] {};

\fill (A) circle (1pt) node[below left] {$A$};
\fill (B) circle (1pt) node[above right] {$B$};
\fill (M) circle (1pt) node[below left] {$M$};
\fill (N) circle (1pt) node[above right] {$N$};
\end{tikzpicture}

\vspace{10pt}

$\overrightarrow{AB} = \overrightarrow{MN}$
\end{minipage}%
\hfill
\begin{minipage}{.24\textwidth}%
\centering
\textbf{4}

\vspace{10pt}

\begin{tikzpicture}[scale=0.8]
\draw[step=1cm, black] (0,0) grid (4,4);

\coordinate (A) at (2,2);
\coordinate (B) at (4,2);
\coordinate (M) at (2,2);
\coordinate (N) at (0,2);

\draw[->, very thick] (A)--(B) node[midway, below] {};
\draw[->, very thick] (M)--(N) node[midway, above] {};

\fill (A) circle (2pt) node[above right] {$A$};
\fill (B) circle (1pt) node[below left] {$B$};
\fill (M) circle (1pt) node[below left] {$M$};
\fill (N) circle (1pt) node[below left] {$N$};
\end{tikzpicture}

\vspace{10pt}

$\overrightarrow{AB} = \overrightarrow{MN}$
\end{minipage}%
\end{center}

\vspace{10pt}

% ------ EXERCICE N°6 -------

{\color{RoyalBlue}\textbf{\underline{Exercice d'application n°6 :}}
\vspace{5pt}
Placer les points pour respecter les égalités.
}

\begin{center}
\begin{minipage}{.24\textwidth}%
\centering
\textbf{1}

\vspace{10pt}

\begin{tikzpicture}[scale=0.7]
\draw[step=1cm, black] (0,0) grid (5,4);

\coordinate (A) at (1,3);
\coordinate (B) at (4,2);
\coordinate (M) at (2,1);
% \coordinate (N) at (3,2);

% \draw[->, very thick] (A)--(B) node[midway, below right] {};
% \draw[->, very thick] (M)--(N) node[midway, below left] {};

\fill (A) circle (2pt) node[below left] {$A$};
\fill (B) circle (2pt) node[above right] {$B$};
\fill (M) circle (2pt) node[below left] {$M$};
% \fill (N) circle (1pt) node[below left] {$N$};
\end{tikzpicture}

\vspace{10pt}

$N$ tel que $\overrightarrow{AB} = \overrightarrow{MN}$
\end{minipage}%
\hfill
\begin{minipage}{.31\textwidth}%
\centering
\textbf{2}

\vspace{10pt}

\begin{tikzpicture}[scale=0.7]
\draw[step=1cm, black] (0,0) grid (6,4);

\coordinate (C) at (1,3);
\coordinate (D) at (3,3);
\coordinate (E) at (5,1);
% \coordinate (N) at (1,3);

% \draw[->, very thick] (A)--(B) node[midway, below] {};
% \draw[->, very thick] (M)--(N) node[midway, above] {};

\fill (C) circle (2pt) node[below left] {$C$};
\fill (D) circle (2pt) node[above right] {$D$};
\fill (E) circle (2pt) node[above right] {$E$};
% \fill (N) circle (1pt) node[below left] {$N$};
\end{tikzpicture}

\vspace{10pt}

$F$ tel que $\overrightarrow{DF} = \overrightarrow{CE}$
\end{minipage}%
\hfill
\begin{minipage}{.17\textwidth}%
\centering
\textbf{3}

\vspace{10pt}

\begin{tikzpicture}[scale=0.7]
\draw[step=1cm, black] (0,0) grid (4,4);

\coordinate (A) at (1,1);
\coordinate (B) at (3,2);
\coordinate (M) at (1,3);
\coordinate (N) at (3,4);

\draw[->, very thick] (A)--(B) node[midway, below] {};
\draw[->, very thick] (M)--(N) node[midway, above] {};

\fill (A) circle (1pt) node[below left] {$A$};
\fill (B) circle (1pt) node[above right] {$B$};
\fill (M) circle (1pt) node[below left] {$M$};
\fill (N) circle (1pt) node[above right] {$N$};
\end{tikzpicture}

\vspace{10pt}

$\overrightarrow{AB} = \overrightarrow{MN}$
\end{minipage}%
\hfill
\begin{minipage}{.24\textwidth}%
\centering
\textbf{4}

\vspace{10pt}

\begin{tikzpicture}[scale=0.7]
\draw[step=1cm, black] (0,0) grid (4,4);

\coordinate (A) at (2,2);
\coordinate (B) at (4,2);
\coordinate (M) at (2,2);
\coordinate (N) at (0,2);

\draw[->, very thick] (A)--(B) node[midway, below] {};
\draw[->, very thick] (M)--(N) node[midway, above] {};

\fill (A) circle (2pt) node[above right] {$A$};
\fill (B) circle (1pt) node[below left] {$B$};
\fill (M) circle (1pt) node[below] {$M$};
\fill (N) circle (1pt) node[below left] {$N$};
\end{tikzpicture}

\vspace{10pt}

$\overrightarrow{AB} = \overrightarrow{MN}$
\end{minipage}%
\end{center}

% -------- Exercice n°7 ---------- %
{\color{RoyalBlue}\textbf{\underline{Exercice d'application n°7 :}}
\vspace{10pt}
\\On considère la figure ci-dessous :
}
\\
\begin{minipage}[t]{.45\textwidth}%
    \vspace{0pt}
    \begin{tikzpicture}[scale=0.8]

    % Grille
    \draw[step=1cm, black, thin] (0,0) grid (10,5);

    % Axes
    % \draw[->] (0,0) -- (10.5,0) node[right] {$x$};
    % \draw[->] (0,0) -- (0,5.5) node[above] {$y$};

    % Points
    \coordinate (A) at (1,2);
    \coordinate (B) at (7,4);
    \coordinate (C) at (5,1);
    \coordinate (D) at (9,1);

    \fill (A) circle (2pt) node[above left] {$A$};
    \fill (B) circle (2pt) node[above left] {$B$};
    \fill (C) circle (2pt) node[below] {$C$};
    \fill (D) circle (2pt) node[below] {$D$};

    % Vecteur v
    \draw[->, ultra thick, blue]
    (3,1) -- (4,4)
    node[midway, right] {$\overrightarrow{v}$};

    \fill (3,1) circle (1pt);
    \fill (4,4) circle (1pt);

    \end{tikzpicture}
\end{minipage}%
\hfill
\begin{minipage}[t]{.55\textwidth}%
    % \vspace{Opt}
    % \\
    \vspace{5pt}
    \begin{enumerate}
        \item Reproduire la figure.
        \item Contruire l'image $E$ du point $A$ par la translation de vecteur $\overrightarrow{CD}$.
        \item Construire le point $F$ tel que $\overrightarrow{BF} = \overrightarrow{CA}$.
        \item Construire le représentant d'origine $A$ du vecteur $\overrightarrow{v}$.
        \item Construire le représentant d'origine $B$ du vecteur $\overrightarrow{v}$.
    \end{enumerate}
\end{minipage}

\end{mdframed}

% ============================================
%   PARTIE 2 : Somme et différence de vecteurs
% ============================================
\newpage
\section*{\textcolor{red}{Partie 2 : Somme et différence de vecteurs}}

\subsection*{\textcolor{red}{2.1 Exemple avec des translations}}

\begin{minipage}[t]{.45\textwidth}%
    Soit {\color{ForestGreen}$t_1$ la translation de vecteur $\overrightarrow{u}$} et {\color{red}
    $t_2$ la translation de vecteur $\overrightarrow{v}$}.
    \vspace{5pt}
    \\Appliquer la translation $t_1$ puis la translation $t_2$ :
    \begin{center}
        $M \quad {\color{ForestGreen} \longmapsto} \quad M_1 \quad {\color{red} \longmapsto} 
        \quad M_2$
    \end{center}
    \vspace{0pt}
    Revient à appliquer la {\color{ForestGreen}translation $t$ de vecteur $\overrightarrow{w}$} :
    \begin{center}
        $\_\_\_ \quad {\color{blue} \longmapsto} \quad \_\_\_$
    \end{center}
    L'enchaînement de deux translations de vecteurs $\color{ForestGreen}\overrightarrow{u}$ et 
    $\color{red}\overrightarrow{v}$ est 
    la translation de vecteur 
    \vspace{10pt}
    \\$\color{blue}\overrightarrow{w} = \_\_\_\_\_\_\_\_\_\_\_\_$.
\end{minipage}%
\hfill
\begin{minipage}[t]{.45\textwidth}%
    \vspace{0pt}
    \begin{tikzpicture}[scale=1]

    % Grille
    % \draw[step=1cm, black, thin] (-1,-1) grid (6,4);

    % Vecteur u
    \draw[->, ultra thick, ForestGreen]
    (0,3) -- (3,4)
    node[midway, above] {$\overrightarrow{u}$};

    % Vecteur v
    \draw[->, ultra thick, red]
    (5,3) -- (6,5)
    node[midway, above left] {$\overrightarrow{v}$};

    \end{tikzpicture}

    \vspace{10pt}

    \begin{tikzpicture}[scale=1]

    % Grille (facultative)
    % \draw[step=1cm, black, thin] (-1,-1) grid (6,4);

    % Points
    \coordinate (A) at (0,0);
    \coordinate (B) at (3,1);
    \coordinate (C) at (4,3);

    % Vecteur u
    \draw[->, ultra thick, ForestGreen]
    (A) -- (B)
    node[midway, below] {$\overrightarrow{u}$};

    % Vecteur v (bout à bout)
    \draw[->, ultra thick, red]
    (B) -- (C)
    node[midway, right] {$\overrightarrow{v}$};

    % Vecteur w = u + v (ferme le triangle)
    \draw[->, ultra thick, blue]
    (A) -- (C)
    node[midway, above left] {$\overrightarrow{w}$};

    % Triangle en pointillés (optionnel mais très lisible)
    \draw[dashed] (A) -- (B) -- (C) -- cycle;

    % Points
    \fill (A) circle (1pt) node[left] {$M$};
    \fill (B) circle (1pt) node[below right] {$M_1$};
    \fill (C) circle (1pt) node[above] {$M_2$};

    \end{tikzpicture}

\end{minipage}%

\vspace{20pt}


% =========================
% Addition de deux vecteurs
% =========================

\subsection*{\textcolor{red}{2.2 Addition de deux vecteurs}}

\vspace{10pt}

\begin{mdframed}[
    linewidth=1pt,
    roundcorner=5pt,
    linecolor=ForestGreen,
    innerbottommargin=10pt
]
{\color{ForestGreen}\textbf{Définition :}}
    \vspace{5pt}
    \\En enchaînant la translation de vecteur $\overrightarrow{u}$ et la translation de vecteur 
    $\overrightarrow{v}$ on obtient une nouvelle translation.
    \vspace{5pt}
    \\Le vecteur associé à cette translation est la somme des vecteurs $\overrightarrow{u}$ et 
    $\overrightarrow{v}$.
    \vspace{5pt}
    \\Il est noté $\bm{\overrightarrow{u} + \overrightarrow{v}}$.
\end{mdframed}

\begin{mdframed}[
    linewidth=1pt,
    roundcorner=5pt,
    linecolor=Red,
    innertopmargin=15pt,
    % innerbottommargin=10pt
]
{\color{red}
\underline{Propriété :} \textbf{Relation de Chasles}
\vspace{5pt}
\\Pour tous points du plan : 
\[
\bm{
\overrightarrow{A\color{blue}B} + \overrightarrow{{\color{blue}B}C} = \overrightarrow{AC}
}
\]
}
\end{mdframed}

\vspace{10pt}

\underline{Remarques :}
\vspace{5pt}
\begin{itemize}
    \item $\overrightarrow{AB} + \overrightarrow{BA} = \overrightarrow{AA} = \overrightarrow{0}$.
    \vspace{5pt}
    \item Soustraire un vecteur c'est \textbf{additionner son } \_\_\_\_\_\_\_\_\_\_\_\_.
    \vspace{5pt}
    \\\underline{Exemple :} En utilisant la relation de Chasles
    \vspace{10pt}
    \\$\overrightarrow{AB} - \overrightarrow{AC} = \_\_\_\_\_\_\_\_\_\_\_$ = $\_\_\_\_\_\_\_\_\_\_\_$ = $\_\_\_\_$
\end{itemize}

\vspace{20pt}

% =========================
% Soustraction de deux vecteurs
% =========================

\subsection*{\textcolor{red}{2.3 Soustraction de deux vecteurs}}

\vspace{5pt}
Pour effectuer la différence des vecteurs $\overrightarrow{u}$ et $\overrightarrow{v}$, on passe à la somme :
\[
    \overrightarrow{u} - \overrightarrow{v} = \overrightarrow{u} + (\overrightarrow{-v})
\]
Pour obtenir la somme des vecteurs $\overrightarrow{u}$ et $-\overrightarrow{v}$, on construit les vecteurs 
$\overrightarrow{u}$ et $\overrightarrow{-v}$ \underline{bout à bout}.

\vspace{10pt}

\begin{center}
    \begin{tikzpicture}[scale=1]

        % Grille
        % \draw[step=1cm, black, thin] (-1,-1) grid (6,4);

        % Vecteur u
        \draw[->, ultra thick, ForestGreen]
        (0,0) -- (2,2)
        node[midway, above left] {$\overrightarrow{u}$};

        % Vecteur v
        \draw[->, ultra thick, red]
        (3,3) -- (6,2)
        node[midway, above] {$\overrightarrow{v}$};

        % Vecteur v
        \draw[->, ultra thick, red]
        (6,0) -- (3,1)
        node[midway, above] {$\overrightarrow{-v}$};

    \end{tikzpicture}

    \vspace{10pt}

    \begin{tikzpicture}[scale=1]
        
        % Grille (facultative)
        % \draw[step=1cm, black, thin] (-1,-1) grid (6,4);

        % Points
        \coordinate (A) at (0,0);
        \coordinate (B) at (2,2);
        \coordinate (C) at (-1,3);

        % Vecteur u
        \draw[->, ultra thick, ForestGreen]
        (A) -- (B)
        node[midway, right] {$\overrightarrow{u}$};

        % Vecteur v (bout à bout)
        \draw[->, ultra thick, red]
        (B) -- (C)
        node[midway, above] {$\overrightarrow{-v}$};

        % Vecteur w = u + v (ferme le triangle)
        \draw[->, ultra thick, blue]
        (A) -- (C)
        node[midway, left] {$\overrightarrow{w}$};

        % Triangle en pointillés (optionnel mais très lisible)
        \draw[dashed] (A) -- (B) -- (C) -- cycle;

        % Points
        \fill (A) circle (0pt) node[left] {};
        \fill (B) circle (0pt) node[below right] {};
        \fill (C) circle (0pt) node[above] {};

    \end{tikzpicture}
\end{center}

\begin{mdframed}[
    linewidth=1pt,
    roundcorner=5pt,
    linecolor=RoyalBlue,
    innertopmargin=10pt,
    innerbottommargin=10pt
]
% ==========================
% Exercice d'application n°8
% ==========================
{\color{RoyalBlue}\textbf{\underline{Exercice d'application n°8 :} Construire une somme de vecteurs}
\vspace{5pt}
}
\\Reproduire les figures suivantes puis placer le point $M$ tel que $\overrightarrow{AM} = \overrightarrow{u}
+ \overrightarrow{v}$.
\\
\begin{minipage}[t]{.24\textwidth}%
    \vspace{0pt}
    \begin{tikzpicture}[scale=0.65]

    % Grille
    \draw[step=1cm, black, thin] (-1,-1) grid (5,4);

    % Points
    \coordinate (A) at (0,0);

    % Vecteur u
    \draw[->, ultra thick, ForestGreen]
    (0,0) -- (2,3)
    node[midway, above left] {$\overrightarrow{u}$};

    % Vecteur v
    \draw[->, ultra thick, red]
    (2,3) -- (4,1)
    node[midway, above right] {$\overrightarrow{v}$};

    \fill (A) circle (2pt) node[left] {$A$};

    \node[below] at (0,3.5) {\textbf{(a)}};

    \end{tikzpicture}
\end{minipage}%
\hfill
\begin{minipage}[t]{.24\textwidth}%
\vspace{0pt}
    \begin{tikzpicture}[scale=0.65]

    % Grille
    \draw[step=1cm, black, thin] (-1,-1) grid (5,4);

    % Points
    \coordinate (A) at (0,2);

    % Vecteur u
    \draw[->, ultra thick, ForestGreen]
    (1,1) -- (1,3)
    node[midway, right] {$\overrightarrow{u}$};

    % Vecteur v
    \draw[->, ultra thick, red]
    (0,0) -- (4,0)
    node[midway, above right] {$\overrightarrow{v}$};

    \fill (A) circle (2pt) node[left] {$A$};

    \node[below] at (0,3.5) {\textbf{(b)}};

    \end{tikzpicture}
\end{minipage}%
\hfill
\begin{minipage}[t]{.24\textwidth}%
\vspace{0pt}
    \begin{tikzpicture}[scale=0.65]

    % Grille
    \draw[step=1cm, black, thin] (-1,-1) grid (5,4);

    % Points
    \coordinate (A) at (2,1);

    % Vecteur u
    \draw[->, ultra thick, ForestGreen]
    (1,0) -- (0,2)
    node[midway, below left] {$\overrightarrow{u}$};

    % Vecteur v
    \draw[->, ultra thick, red]
    (2,2) -- (5,1)
    node[midway, above right] {$\overrightarrow{v}$};

    \fill (A) circle (2pt) node[left] {$A$};

    \node[below] at (0,3.5) {\textbf{(c)}};

    \end{tikzpicture}
\end{minipage}%
\hfill
\begin{minipage}[t]{.24\textwidth}%
\vspace{0pt}
    \begin{tikzpicture}[scale=0.65]

    % Grille
    \draw[step=1cm, black, thin] (-1,-1) grid (5,4);

    % Points
    \coordinate (A) at (1,2);

    % Vecteur u
    \draw[->, ultra thick, ForestGreen]
    (1,0) -- (0,2)
    node[midway, below left] {$\overrightarrow{u}$};

    % Vecteur v
    \draw[->, ultra thick, red]
    (2,1) -- (5,4)
    node[midway, above left] {$\overrightarrow{v}$};

    \fill (A) circle (2pt) node[above] {$A$};

    \node[below] at (0,3.5) {\textbf{(d)}};

    \end{tikzpicture}
\end{minipage}%

\vspace{20pt}


% ==========================
% Exercice d'application n°9
% ==========================

{\color{RoyalBlue}\textbf{\underline{Exercice d'application n°9 :} Travailler avec des vecteurs}
\vspace{5pt}
}
\\$ADCG$ et $AGFE$ sont deux carrés.
\\
\begin{minipage}[t]{.45\textwidth}%
    \vspace{0pt}
    \begin{enumerate}
        \item 
            \begin{enumerate}
                \item Déterminer l'image du point $D$ par la translation de vecteur $\overrightarrow{DA}$ suivie de la 
                tranlsation de vecteur $\overrightarrow{AF}$.
                \item En déduire $\overrightarrow{DA} + \overrightarrow{AF}$.
            \end{enumerate}
        \item Déterminer :
        \begin{enumerate}
            \item $\overrightarrow{DA} + \overrightarrow{AG}$
            \item $\overrightarrow{DA} + \overrightarrow{DG}$
            \item $\overrightarrow{CA} + \overrightarrow{EF}$
            \item $\overrightarrow{EG} + \overrightarrow{DA}$
            \item $\overrightarrow{EG} + \overrightarrow{CA}$
            \item $\overrightarrow{DE} + \overrightarrow{FC}$
        \end{enumerate}
    \end{enumerate}
\end{minipage}%
\hfill
\begin{minipage}[t]{.45\textwidth}%
    \vspace{20pt}
    \begin{tikzpicture}[scale=1]

    % Grille (facultative)
    % \draw[step=1cm, black, thin] (-1,-1) grid (4,6);

    % Points
    \coordinate (A) at (1,3);
    \coordinate (G) at (3,3);
    \coordinate (C) at (3,5);
    \coordinate (D) at (1,5);

    \coordinate (E) at (1,1);
    \coordinate (F) at (3,1);

    % Carré du haut DCGA
    \draw[thick, red] (D) -- (C) -- (G) -- (A) -- cycle;

    % Carré du bas AGFE
    \draw[thick, red] (A) -- (G) -- (F) -- (E) -- cycle;

    % Noms des points
    \node[left]  at (D) {$D$};
    \node[right] at (C) {$C$};
    \node[right] at (G) {$G$};
    \node[left]  at (A) {$A$};
    \node[left]  at (E) {$E$};
    \node[right] at (F) {$F$};

    \end{tikzpicture}

\end{minipage}%

\vspace{20pt}

% ==========================
% Exercice d'application n°10
% ==========================
{\color{RoyalBlue}\textbf{\underline{Exercice d'application n°10 :} Simplifier des égalités de vecteurs}
\vspace{5pt}
}
\\À l'aide de la relation de Chasles, écrire sous forme d'un seul vecteur (si c'est possible !) :
\\
\begin{minipage}[t]{.25\textwidth}%
    \begin{enumerate}
        \item $\overrightarrow{AD} + \overrightarrow{DF}$
        \item $\overrightarrow{CB} + \overrightarrow{CA}$
        \item $\overrightarrow{DF} - \overrightarrow{DG}$
        \item $\overrightarrow{AB} - \overrightarrow{AC}$
    \end{enumerate}
\end{minipage}%
\hfill
\begin{minipage}[t]{.75\textwidth}%
    \begin{enumerate}
        \setcounter{enumi}{4}
        \item $\overrightarrow{RS} + \overrightarrow{AR}$
        \item $\overrightarrow{EG} + \overrightarrow{GT}$
        \item $\overrightarrow{AL} - \overrightarrow{LA}$
        \item $- \overrightarrow{AD} - \overrightarrow{DB}$
    \end{enumerate}
\end{minipage}%
\end{mdframed}

\begin{mdframed}[
    linewidth=1pt,
    roundcorner=5pt,
    linecolor=Red,
    innertopmargin=10pt,
    innerbottommargin=10pt
]
{\color{red}
\underline{Propriété :} \textbf{Caractéristique du parallélogramme}
\vspace{5pt}
\\Dire que $ABCD$ est un parallélogramme revient à dire que $\overrightarrow{AC} = \overrightarrow{AB} + 
\overrightarrow{AD}$.
\begin{center}
    \begin{tikzpicture}[scale=1]

    % Points
    \coordinate (A) at (1,1);
    \coordinate (B) at (4,1);
    \coordinate (D) at (2,3);
    \coordinate (C) at (5,3);

    % Parallélogramme
    \draw[->, very thick] (A) -- (B) node[midway, below] {$\overrightarrow{AB}$};
    \draw[->, very thick] (A) -- (D) node[midway, above left] {$\overrightarrow{AD}$};
    \draw[->, very thick] (A) -- (C) node[midway, above] {$\overrightarrow{AC}$};

    \draw[dashed] (B) -- (C);
    \draw[dashed] (C) -- (D);

    % Points
    \fill (A) circle (1pt) node[below left] {$A$};
    \fill (B) circle (1pt) node[below right] {$B$};
    \fill (C) circle (1pt) node[above right] {$C$};
    \fill (D) circle (1pt) node[above left] {$D$};

    \end{tikzpicture}
\end{center}
}
\end{mdframed}

\underline{Démonstration :}
\vspace{5pt}
\\D'après la relation de Chasles, l'égalité ${\color{red}\overrightarrow{AC}} = \overrightarrow{AB} + \overrightarrow{AD}$ peut 
s'écrire :
\\${\color{red}\overrightarrow{AD} + \overrightarrow{DC}} = \overrightarrow{AB} + \overrightarrow{AD}$
\\Soit $\overrightarrow{DC} = \overrightarrow{AB}$,
\\Soit encore : $ABCD$ est un parallélogramme.

\end{document}
